\section{Diskussion}
\label{sec:diskussion}
 
In diesem Abschnitt werden die Ergebnisse der Knowledge Distillation in den Kontext der beiden vorangegangenen
Meilensteine eingeordnet, die auf demselben Datensatz durchgeführt wurden. In Meilenstein~1 wurden zehn selbst
entworfene CNN-Architekturen von Grund auf trainiert, in Meilenstein~2 wurden vortrainierte ResNet-Modelle
(ResNet-18 bis ResNet-152) mittels Transfer Learning angepasst. Das beste Modell aus Meilenstein~2 diente anschließend
als Teacher für die in dieser Arbeit untersuchte Distillation. Abschließend werden die bereits in
Abschnitt~\ref{sec:quantitative_ergebnisse} beobachteten Effekte der Hyperparameter $\alpha$ und $T$ vertiefend
analysiert.
 
\subsection{Vergleich mit den Architekturen aus Meilenstein~1}
\label{sec:diskussion_m1}
 
Besonders aufschlussreich ist der Vergleich mit Meilenstein~1, da die dort untersuchten Architekturen hinsichtlich der
Parameteranzahl in einer ähnlichen Größenordnung liegen wie das Student-Netz, jedoch durchweg tiefer sind, also über mehr
Schichten verfügen (vgl. Tabelle~\ref{tab:architektur_uebersicht}). Der Vergleich erlaubt daher eine Einschätzung,
inwieweit Knowledge Distillation den Nachteil einer geringeren Netztiefe kompensieren kann.
 
Als direkteste Vergleichsgruppe eignen sich die Architekturen mit Batch-Normalisierung, da auch das Student-Netz diese
Technik einsetzt. Das ohne Teacher trainierte Student-Netz erreicht mit einer Testgenauigkeit von 73{,}16\,\% ein
Ergebnis, das unterhalb aller BatchNorm-Varianten aus Meilenstein~1 liegt: Bereits das einfache \textit{BatchNorm\_Net}
erzielt 79{,}6\,\%, die beste Architektur \textit{BatchNorm\_Dense\_Pool\_Conv\_Net} sogar 84{,}2\,\%. Lediglich
das \textit{BatchNorm\_Dense\_Pool\_Net} liegt mit 73{,}0\,\% auf demselben Niveau wie die Baseline. Dieser Rückstand
ist plausibel auf die geringere Tiefe des Student-Netzes zurückzuführen, die dessen Repräsentationsfähigkeit einschränkt.
 
Durch die Knowledge Distillation verschiebt sich dieses Bild deutlich: Das beste destillierte Modell
($\alpha = 0{,}1$, $T = 4{,}0$) erreicht mit 82{,}11\,\% eine Testgenauigkeit, die nur noch rund 2~Prozentpunkte
unter der besten Architektur aus Meilenstein~1 liegt und drei der fünf BatchNorm-Varianten übertrifft. Das flachere
Student-Netz kommt mit Unterstützung des Teachers also nahezu an das Ergebnis heran, für das in Meilenstein~1 eine
tiefere und aufwendig per Architektursuche optimierte Netzstruktur erforderlich war. Dies deutet darauf hin, dass ein
wesentlicher Teil des Leistungsunterschieds zwischen flachen und tiefen Netzen dieser Größenordnung nicht in der
Kapazität der Netze selbst begründet liegt, sondern in der Qualität des Trainingssignals: Die weichen Zielverteilungen
des Teachers stellen dem Student-Netz zusätzliche Information über die Ähnlichkeitsstruktur der Klassen bereit, die aus
den harten Labels allein nicht hervorgeht.
 
\subsection{Vergleich mit den Transfer-Learning-Modellen aus Meilenstein~2}
\label{sec:diskussion_m2}
 
Die mittels Transfer Learning angepassten ResNet-Modelle aus Meilenstein~2 markieren die obere Leistungsgrenze in dieser
Versuchsreihe: Alle fünf Varianten erreichen Validierungsgenauigkeiten zwischen 96{,}8\,\% und 98{,}9\,\%, wobei das
ResNet-50 mit 98{,}95\,\% am besten abschneidet und daher als Teacher für die Distillation ausgewählt wurde. Auffällig
ist, dass die Leistung innerhalb der ResNet-Familie kaum mit der Netztiefe skaliert was darauf hindeutet, dass die auf ImageNet
vorgelernten Merkmale für den vorliegenden Datensatz bereits weitgehend ausreichen und die Aufgabe für diese Modellklasse nicht
kapazitätslimitiert ist.
 
Gegenüber dem Teacher verbleibt auch nach der Distillation eine erhebliche Lücke von rund 15~Prozentpunkten
(82{,}11\,\% gegenüber 98{,}95\,\%). Das Student-Netz kann das Wissen des Teachers also nur teilweise übernehmen,
was angesichts des großen Kapazitäts- und Tiefenunterschieds sowie des fehlenden Vortrainings zu erwarten war. Die
Distillation ist somit kein Ersatz für Transfer Learning, wenn maximale Genauigkeit das Ziel ist. Sie schließt jedoch
einen relevanten Teil der Lücke zwischen dem frei trainierten Student-Netz und dem Teacher: Von den ursprünglich rund
26~Prozentpunkten Abstand zwischen Baseline und Teacher werden etwa 9~Prozentpunkte durch die Distillation
zurückgewonnen. Für Einsatzszenarien mit beschränkten Ressourcen, etwa auf eingebetteten Systemen, stellt der destillierte Student
daher einen attraktiven Kompromiss zwischen Genauigkeit und Modellgröße dar.
 
\subsection{Einfluss der Hyperparameter}
\label{sec:diskussion_hyperparameter}
 
Die in Abschnitt~\ref{sec:quantitative_ergebnisse} dargestellten Ergebnisse zeigen, dass der Erfolg der Distillation
empfindlich von der Wahl der Hyperparameter $\alpha$ und $T$ abhängt. Über die Konfigurationen hinweg schwankt die
Testgenauigkeit um mehr als 10~Prozentpunkte (71{,}40\,\% bis 82{,}11\,\%), und zwei Kombinationen bleiben sogar
hinter der Baseline ohne Teacher zurück. Eine ungünstige Parameterwahl kann den Nutzen der Distillation demnach nicht
nur aufheben, sondern das Training aktiv verschlechtern.
 
Bei der Temperatur deutet sich ein systematischer Effekt an: Zwei der drei Läufe mit $T = 2{,}0$ zählen zu den
schwächsten Ergebnissen. Eine niedrige Temperatur lässt die weichen Zielverteilungen des Teachers gegen die harten
Labels konvergieren, wodurch der zusätzliche Informationsgehalt über die Klassenähnlichkeiten weitgehend verloren
geht -- der Student erhält dann im Wesentlichen ein redundantes, aber verrauschtes zweites Label-Signal. Die moderate
Temperatur $T = 4{,}0$ erweist sich in Kombination mit $\alpha = 0{,}1$ als optimal, während $T = 6{,}0$ die
Verteilungen bereits so stark glättet, dass die Ergebnisse wieder leicht abfallen.
 
Beim Gewichtungsfaktor zeigt sich, dass die beste Leistung bei der kleinsten untersuchten Gewichtung des
Hard-Label-Anteils ($\alpha = 0{,}1$) erzielt wird, das Signal des Teachers also stark dominieren sollte. Dies ist
konsistent mit der hohen Qualität des Teachers: Bei einer Genauigkeit von annähernd 99\,\% sind dessen Vorhersagen
fast immer korrekt und zugleich informativer als die harten Labels, sodass eine stärkere Gewichtung der
Ground-Truth-Labels kaum zusätzlichen Nutzen stiftet. Ein streng monotoner Zusammenhang lässt sich allerdings weder
für $\alpha$ noch für $T$ ableiten; die Wechselwirkung beider Parameter ist offensichtlich nichtlinear. Da jede
Konfiguration nur zweimal trainiert wurde, ist zudem nicht auszuschließen, dass ein Teil der beobachteten Unterschiede
auf die Zufallsinitialisierung zurückgeht.
